\chapter{Lossless Coding}

\section{Table of Contents}

\begin{itemize}
  \item Core Idea
  \item Key Concepts
  \item Theory and Formulas
  \item Visual Schemes
  \item Examples
\end{itemize}

\section{Core Idea}

\textbf{Lossless coding} maps a sequence of discrete symbols to a bitstream that can be decoded exactly. In multimedia compression it usually comes after prediction, transform, or quantization, and it removes remaining \textbf{statistical redundancy} without adding distortion.

Typical pipeline:

\begin{enumerate}
  \item \textbf{Prediction} reduces sample correlation and produces residuals.
  \item \textbf{Quantization} may map residuals to discrete indexes in lossy systems.
  \item \textbf{Lossless coding} maps indexes to compact bits.
\end{enumerate}

\begin{calloutbox}{red!70!black}{Lossless coding limit}
No lossless code can beat source entropy on average. Good entropy coders try to make average length $\mathcal{L}$ approach entropy $H$ or entropy rate $\mathcal{H}$.
\end{calloutbox}

\section{Key Concepts}

\subsection{Codes and Decodability}

Let the alphabet be:

\[
\mathcal{X}=\{x_1,x_2,\ldots,x_M\}
\]

A code is:

\[
C:x_i\in\mathcal{X}\rightarrow c_i\in\{0,1\}^*
\]

where $c_i$ is a finite binary string and $\ell_i$ is its length.

\begin{center}\fitwidth{%
\begin{tabular}{|l|l|l|}
\hline
\textbf{Code type} & \textbf{Meaning} & \textbf{Practical effect} \\ \hline
\textbf{Fixed-length coding (FLC)} & all symbols use same length $\lceil\log_2 M\rceil$ & simple, robust, ignores probabilities \\ \hline
\textbf{Variable-length coding (VLC)} & probable symbols use shorter codewords & compresses non-uniform sources \\ \hline
\textbf{Prefix / instantaneous code} & no codeword is prefix of another & decodes symbol-by-symbol \\ \hline
\textbf{Uniquely decodable code} & any finite bitstream has one symbol sequence & may require lookahead \\ \hline
\end{tabular}}
\end{center}

FLC rate:

\[
R_{\mathrm{FLC}}=\lceil\log_2 M\rceil
\]

VLC average length:

\[
\mathcal{L}=\sum_i p_i\ell_i
\]

\subsection{Information and Entropy}

\textbf{Self-information} of symbol $x_i$:

\[
I(x_i)=-\log_2 p_i
\]

Rare symbols carry more information; certain symbols carry zero information.

\begin{calloutbox}{red!70!black}{Source entropy}
\[
H(X)=-\sum_i p_i\log_2 p_i
\]

Entropy is the average uncertainty of $X$ and the fundamental lower bound for average lossless code length.
\end{calloutbox}

For a binary source with $P(1)=p$:

\[
H(X)=-p\log_2p-(1-p)\log_2(1-p)
\]

Entropy is maximum at $p=0.5$, where $H=1$ bit/symbol.

\subsection{Joint and Conditional Entropy}

Joint entropy:

\[
H(X,Y)=-\sum_{i,j}p_{i,j}\log_2p_{i,j}
\]

Conditional entropy:

\[
H(X|Y)=\sum_jp_jH(X|Y=y_j)
\]

Chain rule:

\[
H(X,Y)=H(Y)+H(X|Y)=H(X)+H(Y|X)
\]

Important consequences:

\begin{itemize}
  \item If $X$ and $Y$ are independent, $H(X|Y)=H(X)$.
  \item If $X$ is determined by $Y$, $H(X|Y)=0$.
  \item Conditioning cannot increase entropy:
\end{itemize}

\[
H(X|Y)\leq H(X)
\]

\subsection{Maximum Entropy}

For an $M$-symbol alphabet, entropy is maximized by the uniform distribution:

\[
p_i^*=\frac{1}{M}
\]

and:

\[
H_{\max}=\log_2M
\]

\textbf{Meaning for compression:} non-uniformity is compressible; uniform sources are hardest to compress losslessly.

\subsection{Kraft and Prefix Codes}

\begin{calloutbox}{red!70!black}{Kraft inequality}
A prefix code with lengths $\ell_1,\ldots,\ell_M$ exists if and only if:

\[
\sum_{i=1}^{M}2^{-\ell_i}\leq1
\]

If equality holds, the code tree is complete.
\end{calloutbox}

McMillan's theorem says uniquely decodable codes cannot outperform instantaneous prefix codes in average length. Therefore, optimal lossless coding can focus on prefix codes.

\subsection{Shannon Source Coding Theorem}

The ideal relaxed code length is:

\[
\ell_i^*=-\log_2p_i
\]

and the corresponding average length is:

\[
\mathcal{L}^*=\sum_ip_i\ell_i^*=H(X)
\]

\begin{calloutbox}{red!70!black}{Source coding theorem}
Any lossless code satisfies:

\[
\mathcal{L}^*\geq H(X)
\]

Equality is possible only when all symbol probabilities are dyadic:

\[
p_i=2^{-k_i}
\]

With integer lengths $\ell_i=\lceil-\log_2p_i\rceil$:

\[
H(X)\leq\mathcal{L}<H(X)+1
\]
\end{calloutbox}

\subsection{Huffman Coding}

\textbf{Huffman coding} is optimal among prefix codes for a known probability table and a fixed symbol alphabet.

Algorithm:

\begin{enumerate}
  \item Create one leaf per symbol with weight $p_i$.
  \item Merge the two lowest-weight nodes.
  \item Repeat until one root remains.
  \item Assign \texttt{0} and \texttt{1} to branches.
  \item Codewords are root-to-leaf paths.
\end{enumerate}

Strengths:

\begin{itemize}
  \item Optimal for single-symbol prefix coding.
  \item Simple and fast.
  \item Used inside many formats and standards.
\end{itemize}

Limits:

\begin{itemize}
  \item One-symbol Huffman cannot code below 1 bit/symbol for a binary alphabet.
  \item Code lengths are integer, so non-dyadic probabilities create overhead.
  \item Block Huffman improves efficiency but alphabet size grows as $M^K$.
\end{itemize}

\subsection{Block Coding}

Instead of coding single symbols, code blocks:

\[
X^K=(X_1,\ldots,X_K)
\]

Block entropy obeys:

\[
H(X_1,\ldots,X_K)\leq\sum_{i=1}^{K}H(X_i)
\]

Huffman bound for a block:

\[
\mathcal{L}<H(X^K)+1
\]

Per original symbol:

\[
\mathcal{L}_S<\frac{H(X^K)}{K}+\frac{1}{K}
\]

\textbf{Meaning:} the Huffman overhead is spread over $K$ symbols, but probability estimation and codebook size become expensive.

\subsection{Arithmetic Coding}

\textbf{Arithmetic coding} encodes an entire sequence as a subinterval of $[0,1)$. Each symbol narrows the current interval according to its probability.

For a sequence $x^n$:

\[
P(x^n)=\prod_{i=1}^{n}p(x_i)
\]

The required number of bits is close to:

\[
-\log_2P(x^n)
\]

Average per-symbol length:

\[
\mathcal{L}<H(X)+\frac{2}{n}
\]

So:

\[
\mathcal{L}\rightarrow H(X)
\]

as $n\rightarrow\infty$.

Arithmetic coding solves the main Huffman problems:

\begin{itemize}
  \item Works efficiently even when entropy is below 1 bit/symbol.
  \item Avoids exponential block-code alphabets.
  \item Handles adaptive and context-dependent probabilities.
\end{itemize}

\subsection{Context-Based and Adaptive Coding}

\textbf{Adaptive coding} updates probabilities during encoding and decoding using the same rule on both sides.

\textbf{Context coding} estimates:

\[
P(X_n| \text{past context})
\]

instead of a single unconditional $P(X_n)$. If context captures the significant past, conditional entropy drops:

\[
H(X_n|\text{context}) < H(X_n)
\]

For alphabet size $M$ and context length $N_S$:

\[
N_C=M^{N_S}
\]

possible contexts exist. Too many contexts cause sparse training data; too few miss dependencies.

\subsection{Exponential-Golomb Coding}

\textbf{Exp-Golomb} is a universal integer code. It needs no probability table and is useful for metadata and sparse residuals.

Unsigned code for $n\geq0$:

\begin{enumerate}
  \item If $n=0$, code is \texttt{1}.
  \item If $n\geq1$, write $n+1$ in binary with $b=\lfloor\log_2(n+1)\rfloor+1$ bits.
  \item Prefix it with $b-1$ zeros.
\end{enumerate}

Signed mapping:

\[
m(n)=
\begin{cases}
2n-1 & n>0\\
-2n & n\leq0
\end{cases}
\]

Then encode $m$ as unsigned.

\subsection{Dictionary Coding and LZW}

Dictionary coding learns repeated patterns from the sequence itself. \textbf{LZW} initializes the dictionary with one-symbol entries, then adds longer patterns as they are encountered.

Encoding idea:

\begin{enumerate}
  \item Find longest prefix $W$ already in dictionary.
  \item If $W+K$ is new, output index of $W$.
  \item Add $W+K$ to dictionary.
  \item Continue from $K$.
\end{enumerate}

Decoder builds the same dictionary because updates are deterministic. Dictionary does not need to be transmitted.

\begin{calloutbox}{red!70!black}{Dictionary asymptotic optimality}
For stationary ergodic sources, dictionary-based coding is asymptotically optimal:

\[
L_n\xrightarrow{n\to\infty}\mathcal{H}(X)
\]
\end{calloutbox}

\subsection{Standards and Practical Codecs}

\begin{center}\fitwidth{%
\begin{tabular}{|l|l|l|}
\hline
\textbf{Standard / codec} & \textbf{Main tools} & \textbf{Best use} \\ \hline
\textbf{JBIG-1} & context arithmetic coding & bi-level images, progressive coding \\ \hline
\textbf{JBIG-2} & segmentation, symbol dictionaries, arithmetic coding & text/halftone documents, PDF \\ \hline
\textbf{JPEG-LS} & MED prediction, context modeling, Golomb-Rice & natural photos, medical images \\ \hline
\textbf{PNG} & prediction, LZ77/Deflate, Huffman & graphics, icons, text-heavy images \\ \hline
\end{tabular}}
\end{center}

JPEG-LS \textbf{MED predictor} uses causal neighbors:

\begin{center}\fitwidth{%
\begin{tabular}{|l|l|}
\hline
\textbf{C} & \textbf{B} \\ \hline
A & X \\ \hline
\end{tabular}}
\end{center}

\[
\hat{x}=
\begin{cases}
\min(A,B) & C\geq\max(A,B)\\
\max(A,B) & C\leq\min(A,B)\\
A+B-C & \text{otherwise}
\end{cases}
\]

This avoids overshoot around sharp edges.

\subsection{Neural Lossless Coding}

Neural lossless coding uses a neural network to estimate a probability model $Q$ for arithmetic coding.

Training minimizes cross-entropy:

\[
H(P,Q)=H(P)+D_{KL}(P\parallel Q)
\]

where:

\begin{itemize}
  \item $H(P)$ is the irreducible source entropy.
  \item $D_{KL}(P\parallel Q)$ is redundancy from model mismatch.
\end{itemize}

Autoregressive models use:

\[
P(x_1,\ldots,x_n)=\prod_{i=1}^{n}P(x_i|x_1,\ldots,x_{i-1})
\]

They can model long-range dependencies, but decoding may require one model evaluation per pixel.

Neural predictive coding replaces a linear predictor with:

\[
\hat{x}_i=f_{NN}(\text{local context})
\]

and codes residuals. For \texttt{house.pgm}, source reports:

\begin{itemize}
  \item 2D linear predictor residual entropy: \textbf{2.830 bpp}.
  \item MLP predictor residual entropy: \textbf{2.704 bpp}.
\end{itemize}

Total rate must include model cost:

\[
R_{\text{total}}=R_{\text{residuals}}+\frac{L_{\text{model}}}{N}
\]

Limitations:

\begin{itemize}
  \item high decoding complexity,
  \item hardware requirements,
  \item generalization gap,
  \item bit-exact determinism issues.
\end{itemize}

\section{Theory and Formulas}

\subsection{Entropy Bound and Optimal Lengths}

Given probabilities $p_i$, ideal lengths solve:

\[
\min_{\ell_i}\sum_ip_i\ell_i
\quad\text{s.t.}\quad
\sum_i2^{-\ell_i}=1
\]

Using a Lagrangian:

\[
J(\ell)=\sum_ip_i\ell_i+\lambda\left(\sum_i2^{-\ell_i}-1\right)
\]

Relaxing integer constraints:

\[
\ell_i^*=-\log_2p_i
\]

and:

\[
\mathcal{L}^*=H(X)
\]

This explains why optimal entropy coding assigns short codewords to probable symbols and long codewords to rare symbols.

\subsection{Kraft Tree Interpretation}

In a binary tree of maximum depth $L_{\max}$, a codeword of length $\ell_i$ blocks:

\[
2^{L_{\max}-\ell_i}
\]

leaves. Prefix codewords block disjoint leaf sets, therefore:

\[
\sum_i2^{L_{\max}-\ell_i}\leq2^{L_{\max}}
\]

Dividing by $2^{L_{\max}}$ gives Kraft:

\[
\sum_i2^{-\ell_i}\leq1
\]

\subsection{Conditioning and Context Coding}

For sources with memory:

\[
\mathcal{H}=\lim_{n\to\infty}H(X_n|X_{n-1},\ldots,X_1)
\]

Context coding approximates this entropy rate by conditioning on a finite significant past.

For a binary image with previous pixel as context:

\[
H(X|\square)=0.322,\qquad H(X|\blacksquare)=0.918
\]

Weighted average:

\[
H(X|Y)=0.406\ \text{bpp}
\]

This is below single-pixel entropy $H(X)=0.586$ bpp because neighboring pixels are correlated.

\subsection{Huffman vs. Arithmetic Overhead}

Huffman on single symbols:

\[
H(X)\leq\mathcal{L}<H(X)+1
\]

Huffman on $K$-symbol blocks:

\[
\frac{\mathcal{L}}{K}<\frac{H(X^K)}{K}+\frac{1}{K}
\]

Arithmetic coding on length-$n$ sequences:

\[
\mathcal{L}<H(X)+\frac{2}{n}
\]

Arithmetic coding therefore approaches entropy without exponential block alphabets.

\subsection{House Image Result}

For \texttt{house.pgm}, original bit depth is 8 bpp and direct entropy is:

\[
H=7.056\ \text{bpp}
\]

\begin{center}\fitwidth{%
\begin{tabular}{|l|l|l|l|l|}
\hline
\textbf{Coding input} & \textbf{Entropy} & \textbf{Huffman} & \textbf{Exp-Golomb} & \textbf{ZIP} \\ \hline
Direct pixels & 7.06 & 7.08 & 11.32 & 4.00 \\ \hline
1D residuals & 3.31 & 3.38 & 3.43 & 3.23 \\ \hline
2D residuals & 2.83 & 2.89 & 2.94 & 3.14 \\ \hline
\end{tabular}}
\end{center}

Main result:

\[
7.06\ \text{bpp}\rightarrow2.83\ \text{bpp}
\]

through 2D prediction before entropy coding.

\section{Visual Schemes}

\subsection{Entropy Coding Pipeline}

\begin{center}\includegraphics[width=\linewidth]{ch03/entropy-coding-pipeline.png}\end{center}
Lossless coding sits after prediction and quantization/index generation, where it removes statistical redundancy from symbol indexes.

\begin{center}\footnotesize\itshape Diagram (Mermaid source):\end{center}
\begin{verbatim}
flowchart LR
    A[Samples] --> B[Prediction]
    B --> C[Residuals]
    C --> D[Quantized indexes]
    D --> E[Entropy coder]
    E --> F[Bitstream]
    F --> G[Entropy decoder]
    G --> H[Indexes]
    H --> I[Reconstruction path]
\end{verbatim}

\subsection{Entropy and Prefix Codes}

\begin{center}\includegraphics[width=0.833\linewidth]{ch03/binary-entropy-curve.png}\end{center}
Binary entropy is maximum at equal probabilities and decreases when one symbol becomes predictable.

\begin{center}\includegraphics[width=\linewidth]{ch03/kraft-tree-proof.png}\end{center}
Kraft inequality follows from prefix codewords occupying disjoint subtrees in a binary tree.

\begin{center}\includegraphics[width=0.75\linewidth]{ch03/huffman-tree-example.png}\end{center}
Huffman coding builds a prefix tree by repeatedly merging the least probable symbols.

\subsection{Huffman and Predictive Coding}

\begin{center}\includegraphics[width=0.833\linewidth]{ch03/binary-image-huffman-example.png}\end{center}
Single-symbol Huffman on a binary image cannot go below 1 bpp even when entropy is below 1 bpp.

\begin{center}\includegraphics[width=\linewidth]{ch03/predictive-huffman-gain.png}\end{center}
Predictive quantization plus Huffman coding improves rate-distortion performance by exploiting residual probability imbalance.

\subsection{Arithmetic Coding}

\begin{center}\includegraphics[width=\linewidth]{ch03/arithmetic-coding-interval.png}\end{center}
Arithmetic coding refines an interval for the whole sequence, making it behave like efficient high-order block coding.

\begin{center}\footnotesize\itshape Diagram (Mermaid source):\end{center}
\begin{verbatim}
flowchart LR
    A[Initial interval 0-1] --> B[Symbol probability model]
    B --> C[Refine interval]
    C --> D[Next symbol]
    D --> B
    C --> E[Final interval]
    E --> F[Binary tag]
\end{verbatim}

\subsection{LZW Dictionary Coding}

\begin{center}\includegraphics[width=\linewidth]{ch03/lzw-encoding-example.png}\end{center}
LZW encoding emits dictionary indexes while adding newly observed repeated patterns.

\begin{center}\includegraphics[width=\linewidth]{ch03/lzw-decoding-example.png}\end{center}
LZW decoding reconstructs the same dictionary deterministically, so the dictionary is not transmitted.

\begin{center}\footnotesize\itshape Diagram (Mermaid source):\end{center}
\begin{verbatim}
flowchart LR
    A[Input stream] --> B[Find longest dictionary prefix]
    B --> C[Emit prefix index]
    C --> D[Add prefix plus next symbol]
    D --> B
    C --> E[Index stream]
\end{verbatim}

\subsection{Prediction Results on House Image}

\begin{center}\includegraphics[width=\linewidth]{ch03/house-direct-coding.png}\end{center}
Direct pixel coding remains close to 7 bpp entropy; dictionary coding can help because repeated byte patterns exist.

\begin{center}\includegraphics[width=0.75\linewidth]{ch03/house-1d-prediction.png}\end{center}
1D DPCM prediction sharply reduces residual entropy compared with direct pixel coding.

\begin{center}\includegraphics[width=0.75\linewidth]{ch03/house-2d-prediction.png}\end{center}
2D adaptive prediction reduces entropy further by exploiting vertical and horizontal image structure.

\section{Examples}

\begin{calloutbox}{green!55!black}{Code decodability}
\begin{center}\fitwidth{%
\begin{tabular}{|l|l|l|l|l|}
\hline
\textbf{Symbol} & \textbf{Code 1} & \textbf{Code 2} & \textbf{Code 3} & \textbf{Code 4} \\ \hline
A & \texttt{0} & \texttt{0} & \texttt{0} & \texttt{0} \\ \hline
B & \texttt{0} & \texttt{1} & \texttt{10} & \texttt{01} \\ \hline
C & \texttt{1} & \texttt{00} & \texttt{110} & \texttt{011} \\ \hline
D & \texttt{10} & \texttt{11} & \texttt{111} & \texttt{0111} \\ \hline
\end{tabular}}
\end{center}

\textbf{Takeaway:} Code 3 is prefix and instantaneously decodable. Code 4 is decodable with delay. Codes 1 and 2 are not uniquely decodable.
\end{calloutbox}

\begin{calloutbox}{green!55!black}{French text entropy}
Alphabet size $M=26$.

\[
R_{\mathrm{FLC}}=\lceil\log_2 26\rceil=5\ \text{bits/symbol}
\]

Source entropy:

\[
H=3.999\ \text{bits/symbol}
\]

Any VLC satisfies $\mathcal{L}\geq3.999$. Huffman gives about $4.041$ bits/symbol, so compression ratio is about $5/4.041=1.238$.
\end{calloutbox}

\begin{calloutbox}{green!55!black}{Huffman on six symbols}
\begin{center}\fitwidth{%
\begin{tabular}{|l|l|l|l|}
\hline
\textbf{Symbol} & \textbf{Probability} & \textbf{Codeword} & \textbf{Length} \\ \hline
A & 0.40 & \texttt{0} & 1 \\ \hline
B & 0.20 & \texttt{100} & 3 \\ \hline
C & 0.15 & \texttt{101} & 3 \\ \hline
D & 0.15 & \texttt{110} & 3 \\ \hline
E & 0.05 & \texttt{1110} & 4 \\ \hline
F & 0.05 & \texttt{1111} & 4 \\ \hline
\end{tabular}}
\end{center}

Average length:

\[
\mathcal{L}=2.3\ \text{bits/symbol}
\]

Entropy:

\[
H\approx2.2464\ \text{bits/symbol}
\]

Huffman is close to the entropy bound.
\end{calloutbox}

\begin{calloutbox}{green!55!black}{Binary image and Huffman block coding}
For a black-white image:

\[
P(\square)=86.7\%,\qquad P(\blacksquare)=13.3\%
\]

Single-pixel entropy:

\[
H(X)=0.586\ \text{bpp}
\]

Single-symbol Huffman still needs:

\[
\mathcal{L}=1\ \text{bpp}
\]

because a binary alphabet needs at least one bit per symbol.

Two-pixel block coding gives:

\[
H([X_1X_2])/2=0.511\ \text{bpp},\qquad
\mathcal{L}_S=0.65\ \text{bpp}
\]

Four-pixel block coding gives:

\[
H([X_1X_2X_3X_4])/4=0.383\ \text{bpp},\qquad
\mathcal{L}_S=0.433\ \text{bpp}
\]

\textbf{Takeaway:} larger blocks move Huffman closer to entropy but increase alphabet complexity.
\end{calloutbox}

\begin{calloutbox}{green!55!black}{Arithmetic coding for sequence ACFD}
Probabilities: A=0.4, B=0.2, C=0.15, D=0.15, E=0.05, F=0.05.

Interval refinement:

\begin{center}\fitwidth{%
\begin{tabular}{|l|l|l|}
\hline
\textbf{Step} & \textbf{Symbol} & \textbf{Interval} \\ \hline
start & - & $[0,1)$ \\ \hline
1 & A & $[0,0.4)$ \\ \hline
2 & C & $[0.24,0.30)$ \\ \hline
3 & F & $[0.297,0.30)$ \\ \hline
4 & D & $[0.29925,0.2997)$ \\ \hline
\end{tabular}}
\end{center}

Any binary number inside final interval identifies the whole sequence.
\end{calloutbox}

\begin{calloutbox}{green!55!black}{Exp-Golomb codes}
\begin{center}\fitwidth{%
\begin{tabular}{|l|l|l|}
\hline
\textbf{$n$} & \textbf{$n+1$ binary} & \textbf{unsigned Exp-Golomb} \\ \hline
0 & \texttt{1} & \texttt{1} \\ \hline
1 & \texttt{10} & \texttt{010} \\ \hline
2 & \texttt{11} & \texttt{011} \\ \hline
3 & \texttt{100} & \texttt{00100} \\ \hline
7 & \texttt{1000} & \texttt{0001000} \\ \hline
\end{tabular}}
\end{center}

\textbf{Takeaway:} efficient for small non-negative integers and simple to decode.
\end{calloutbox}

\begin{calloutbox}{green!55!black}{LZW binary sequence}
Input starts:

\begin{verbatim}
0 0 0 1 0 0 0 0 0 0 1 0 ...
\end{verbatim}

Initial dictionary:

\begin{center}\fitwidth{%
\begin{tabular}{|l|l|}
\hline
\textbf{Code} & \textbf{Pattern} \\ \hline
\texttt{0} & \texttt{0} \\ \hline
\texttt{1} & \texttt{1} \\ \hline
\end{tabular}}
\end{center}

New entries include \texttt{00}, \texttt{001}, \texttt{10}, \texttt{000}, and longer repeated patterns. Emitted values are dictionary indexes, not final raw bits.

\textbf{Takeaway:} repeated patterns become single indexes; decoder reconstructs the same dictionary.
\end{calloutbox}

\begin{calloutbox}{green!55!black}{JPEG-LS vs PNG}
\begin{center}\fitwidth{%
\begin{tabular}{|l|l|l|}
\hline
\textbf{Codec} & \textbf{Core method} & \textbf{Best source} \\ \hline
JPEG-LS & prediction + context + Golomb-Rice & natural/medical images \\ \hline
PNG & LZ77/Deflate + Huffman & graphics, icons, text \\ \hline
\end{tabular}}
\end{center}

\textbf{Takeaway:} both are lossless, but their models target different source statistics.
\end{calloutbox}

\begin{calloutbox}{green!55!black}{Neural model cost}
Small MLP with about 100 weights at 32 bits each:

\[
L_{\text{model}}=3200\ \text{bits}
\]

For a $512\times512$ image and residual rate $2.704$ bpp:

\[
R=2.704+\frac{3200}{512^2}\approx2.716\ \text{bpp}
\]

\textbf{Takeaway:} neural coding must account for model transmission unless the model is shared.
\end{calloutbox}

\begin{calloutbox}{green!55!black}{Method selection}
\begin{center}\fitwidth{%
\begin{tabular}{|l|l|l|}
\hline
\textbf{Source characteristic} & \textbf{Recommended method} & \textbf{Reason} \\ \hline
Memoryless & Huffman or arithmetic & matches entropy bound \\ \hline
Stationary with memory & context arithmetic & exploits conditional probabilities \\ \hline
Locally stationary & adaptive arithmetic & tracks changing statistics \\ \hline
Repeating strings & dictionary coding & learns recurring patterns \\ \hline
Complex long-range dependencies & neural models & learns nonlinear probability models \\ \hline
\end{tabular}}
\end{center}

\textbf{Rule of thumb:} Huffman/Exp-Golomb for speed, arithmetic/context for high compression, dictionary for universal files, neural models for best compression when complexity is acceptable.
\end{calloutbox}
